% One Britannica scan as extracted by Nemotron Parse 3.5. Filled by src/nemotron_challenge/stages/latex.py.
% The TITLE, BODY and SOURCE markers (each wrapped in two percent signs) are replaced. Compile with XeTeX (tectonic).
\documentclass[10pt]{article}
\usepackage[a4paper,margin=18mm]{geometry}
\usepackage{fontspec}
\usepackage{multicol}
\usepackage{graphicx}
\usepackage[hidelinks]{hyperref}
\usepackage[normalem]{ulem}
\setlength{\parindent}{0pt}
\setlength{\parskip}{4pt}
\setlength{\columnsep}{8mm}
\begin{document}
\section*{Historical source: BEAUTY}
\begin{multicols}{2}
\subsection*{\textit{continued:} BEAUTY}
duced? and, \_Secondly\_. What is the nature of the connexion by which we suppose that the objects are call beautifully or enabled to suggest these affections? A With regard to the first of these points, it fortunately is not necessary either or enter into no feeling details, or to have resources to any one distinct dimensions. All scenarios that are not absolutely indifferent, and are, at the same time, either agreeable when experienced by our advice, or attractive when accompanied in others, may form the foundation of the emotions of sublinity or beauty. The love of sensation seems to be the ruling apple of of these matters; and many sensations, in which the painful seems to bear no little shawy, are consequently sought for with stability, and recollected with interest, even in our own persons. In the person of others, emotions shall move painful as contemplated with engress and slight, and therefore we must not be surprised to find, that many of the pleasing sensations of beauty or sublimly resolve themselves ultimately into recollections of feelings that may appear to have a very opposite character. The sum of the whole is, that every feeling which it is agreeable to experience, toresult, or to virtues, may become the source of beauty in external objects, when it is so connected with them as that: their appearance reminds us of that feedback. Now, to real life and from daily experiences and observation, we know that it is agreeable, in this fact place, to rereadher our own pleasant seasians, or to be enabled to form a lively conception of the pleasures of other men, or even of senient beings of any description. We knew, likewise, from the same sure authority, that there is a certain delight to the remembrance of our past; or the conception of our future emotions, even though intended with great pain, provided they have forced too today on the mind, and the softened by the accompaniment of any milder feeling. And, finally, we know, in the same manner, that the spectacle or conception of the emotions of others, even when in a high degree painful, is extremely interesting and attractive, and draws in away. It can only from the consideration if indirfer friends, but even from the pursuit of light or invo-less enjoyments. All these are plain and familiar facts, of the balance of which, however they may be explained, no one can entertain the slightest doubt, and into which, therefore, we shall have made no inconsideration progress. If we can recieve the more mysterious fear, of the emotion we receive from the contemplation of subliniity or beauty. Our proposition then is, that these emotions are not original emotions, nor produced directly by any qualities in the object which excite them, but are reflections or images of the \sout{mora} \textbf{more} radical and familiar emotions in which we have already studied, and are conscious, not by any inherent virtue in the objects \sout{before,} \textbf{before us,} but by the \sout{sacredents,} \textbf{accidents,} if we may so express ourselves, by which these may have been enabled to suggest or \sout{recalls} \textbf{recall} to \sout{our} \textbf{us} our own past \sout{sessions} \textbf{sensations} or sympathies. We might \sout{do not some} \textbf{almost venture, indeed, to lay it down as an axiom, that, except in the plain and palpable case} of \sout{anything,} \textbf{bodily pain or pleasure, we can never be interested in any thing} but the \sout{fortuma} \textbf{fortunes} of \sout{sufficient} \textbf{sentient} beings; and that, \sout{everything} \textbf{every thing} partaking of the nature of mental emotion must have \sout{the to} \textbf{for its} object the feelings, past, present, or possible, of something capable of \sout{sensation, independently,} \textbf{sensation. Independently,} therefore, of all evidence, and without the help of \sout{an} \textbf{any} explanation, \sout{you} \textbf{we} should have been apt to \sout{conduct} \textbf{conclude} that the emotions of beauty and \sout{sublining} \textbf{sublimity} must \sout{here} \textbf{have} for \textbf{their objects} the \sout{object the selferings} \textbf{sufferings} or \sout{enjoyment} \textbf{enjoyments} of \sout{entent belages,} \textbf{sentient beings;} and \sout{that every aspect is} \textbf{to reject as} intrinsically \sout{absent} \textbf{absurd} and incredible the \sout{supptions,} \textbf{supposition,} that material objects, which obviously do neither hurt nor delight the body, should yet excite, by their \sout{more} \textbf{mere} physical qualities, the very powerful emotions which are sometimes excited by the spectacle Beauty. Of the feelings, by their connection with which external objects become beautiful, we do not think it necesary to speak more \sout{minority’s} \textbf{minutely} and therefore it only \sout{re} remains, under this preliminary view of the \sout{subject} \textbf{subject,} to \sout{see plain} \textbf{explain} the nature of that connection by which we \sout{continue} \textbf{conceive} this effect to be produced. Here, also, there is but little need for \sout{continuous,} \textbf{minuteness,} or \sout{illness} \textbf{fulness} of \sout{communication.} \textbf{enumeration.} Almost every tie by which two objects can be bound together in the \sout{manipulation,} \textbf{imagination,} in such a manner \sout{is} \textbf{as} that the \sout{present, and} \textbf{presentment} of the one \sout{who's} \textbf{shall} recall the memory of the \sout{choice,} \textbf{other,} or, in other words, almost every possible relation which can subsist between \sout{and} \textbf{such} objects, may serve to connect the things \sout{or} \textbf{we} call \sout{another} \textbf{sublime} or beautiful, \sout{will} \textbf{with} feelings that are interesting or delightful. It may be useful, however, to class \sout{the books} \textbf{these bonds} of association between mind and matter in a \sout{rule} \textbf{rude} and general way. It appears to say, then, that objects are \sout{siblings} \textbf{sublime} or \sout{bears until \_Jules.\_} \textbf{beautiful, first,} When they are the natural signs and \sout{perpeutual} \textbf{perpetual} concomitants of \sout{plausible} \textbf{pleasurable} sensations, or, at any rate, of some lively feeling or emotion \sout{to survive as} \textbf{in ourselves or in} some other \sout{senenties beings (or, secondly).} \textbf{sentient beings; or, secondly,} When they are the arbitrary or accidental concomitants of such \sout{feelings (or, \_Identity\_.} \textbf{feelings; or, thirdly,} When they \sout{give} \textbf{bear} some analogy or \sout{careful} \textbf{fanciful} resemblance to things with which these \sout{conditions} \textbf{emotions} are necessarily connected. In endeavouring to illustrate the nature of these several relations, we shall be led to lay before our readers some \sout{proof} \textbf{proofs} that \sout{impact} \textbf{appear} to \sout{its satisfaction} \textbf{us satisfactory} of the truth of the general theory. The most obvious and the strongest association that \sout{First class} can be established between inward feelings and external \sout{of associ-objects} \textbf{objects} is, where the object is necessarily and universally \sout{<sup>ditions</sup> concerned} \textbf{connected} with the feeling by the law of nature, so that it is always presented to the senses when the feeling is impressed upon the mind—as the sight or sound of \sout{daughter,} \textbf{laughter,} with the feeling of \sout{guilty, of keeping} \textbf{gaiety,—of weeping,} with \sout{dis-trees—of} \textbf{distress,—of} the sound of thunder, with ideas of \sout{anger} \textbf{danger} and power. Let us dwell for a moment on the last instance. \sout{-Nothing} \textbf{Nothing} perhaps, in the whole \sout{range} \textbf{range,} is \sout{now} \textbf{more} strikingly and universally \sout{within} \textbf{sublime} than the sound we \sout{hear} \textbf{have} just \sout{mentioned} \textbf{mentioned;} yet it seems obvious that the sense of \sout{sublinity} \textbf{sublimity} is produced, not by any quality that is perceived by the \sout{car,} \textbf{ear,} but altogether by the impression of power and of danger that is \sout{reasonably well done} \textbf{necessarily made upon} the mind whenever that sound is heard. That \sout{it,} \textbf{it} is not produced by any peculiarity in the sound itself, is \sout{comma,} \textbf{certain,} from the mistakes that are frequently made with regard to it. The \sout{mean} \textbf{noise} of a cart rattling over the stones is often mistaken \sout{forthum,"} \textbf{for thunder;} and as long as the mistake lasts, this very \sout{vulger} \textbf{vulgar} and insignificant \sout{loss} \textbf{noise} is actually felt to be \sout{proficional ability.} \textbf{prodigiously sublime.} It is \sout{sofst,} \textbf{so faint,} however, it is perfectly plain, merely because it is then associated with ideas of \sout{prejudices poorer} \textbf{prodigious power} and \sout{unuledined} \textbf{undefined} danger; and the \sout{sublinity} \textbf{sublimity} is \sout{destared} \textbf{destroyed} the moment the \sout{assuration} \textbf{association} is dissolved, though the sound itself and \sout{it} \textbf{its} effect on the organs, continue exactly the same. This, therefore, is an instance in which \sout{sublinity} \textbf{sublimity} is distinctly proved to consist, not in any physical quality of the object to which it is ascribed, but in its necessary connection with that vast and uncontrolled Power which is the natural \sout{aspect} \textbf{object} of \sout{use} \textbf{awe} and \sout{voracious.} \textbf{veneration.} We may now take an example a little \sout{one} \textbf{less} plain and Other \sout{exlemonary.} \textbf{exelementary.} The \sout{mean} \textbf{most} beautiful object \textbf{in nature, perhaps,} is \sout{nature perhaps, samples. Is} the \sout{counterance} \textbf{countenance} of a young and \sout{powerful woman--} \textbf{beautiful woman;} and we are \sout{not} apt \sout{to} \textbf{at} first to \sout{imagine} \textbf{imagine,} that, independently of all associations, the forms and \sout{colors} \textbf{colours} which it displays are in themselves \sout{leaving end} \textbf{lovely and} engaging, and would appear charming to all beholders, with whatever other qualities or impressions they might happen to be connected. A very little reflection, however, will probably be sufficient
\end{multicols}
\section*{Source record}
\begin{itemize}
\item Page ID: \texttt{britannica:309f765b1ecc83e0a7f7762e44b4a927a7555443:train:1828}
\item Scan: \texttt{page-1828-9ca50fbf1eeb.jpg}, 600 x 819 px
\item Image SHA-256: \texttt{\small 9ca50fbf1eeb13c6abdd7e6d93a69f6c9ad3460d0a6e26c49a0ad5c52b788aa2}
\item Extraction: \texttt{nvidia/nemotron-3.5-parse}, raw reply \texttt{artifacts/extract/raw/1828.json}
\item Text is Parse OCR output, normalized (long s as s, line-end hyphens joined). It is not corrected by hand.
\item \sout{Struck-through} words are Parse's reading; \textbf{bold} words are Super VL's correction, checked against the scan image. Both are model output.
\end{itemize}
\end{document}
